\documentclass[11pt]{article}

\usepackage[a4paper,margin=1in]{geometry}
\usepackage[T1]{fontenc}
\usepackage[utf8]{inputenc}
\usepackage{amsmath,amssymb}
\usepackage{booktabs}
\usepackage{longtable}
\usepackage{hyperref}
\usepackage{enumitem}

\hypersetup{
  colorlinks=true,
  linkcolor=blue,
  urlcolor=blue,
  citecolor=blue
}

\title{Unitary Synthesis Methodology in the Current Codebase}
\author{}
\date{\today}

\begin{document}
\maketitle
\tableofcontents

\section{Introduction}

This document describes the methodology currently implemented in this repository for quantum unitary synthesis. It is grounded in the active code paths in \texttt{main.py}, \texttt{gqe.py}, \texttt{pipeline.py}, \texttt{model.py}, \texttt{cost.py}, \texttt{continuous\_optimizer.py}, \texttt{pareto.py}, \texttt{pareto\_gd\_optimizer.py}, \texttt{operator\_pool.py}, \texttt{target.py}, and \texttt{config.yml}.

The repository solves unitary synthesis as a hybrid discrete-continuous optimization problem:

\begin{itemize}[leftmargin=*]
\item the discrete part is the circuit structure, represented as a sequence of gate tokens sampled by a transformer policy;
\item the continuous part is the optimization of rotation angles for parameterized gates;
\item the training signal is fidelity-driven, optionally augmented by Pareto-style pressure toward lower depth and fewer CNOT gates.
\end{itemize}

\section{Formal Problem Definition}

Let:

\begin{itemize}[leftmargin=*]
\item $n$ be the number of qubits,
\item $d = 2^n$ be the Hilbert-space dimension,
\item $U_{\mathrm{target}} \in U(d)$ be the target unitary,
\item $V$ be the discrete gate vocabulary induced by the operator pool,
\item $N = \texttt{max\_gates\_count}$ be the fixed sequence length used by the current model.
\end{itemize}

In the current code, a candidate circuit is represented by a token sequence
\[
x = (x_1, \ldots, x_N), \qquad x_t \in V
\]
and, for parameterized gates, a continuous parameter vector
\[
\theta = (\theta_1, \ldots, \theta_m).
\]

Each token $x_t$ maps to a gate matrix $G_t$. For non-parameterized gates, $G_t$ is fixed. For parameterized gates, $G_t$ depends on the relevant entry of $\theta$. The synthesized circuit unitary is
\[
U(x,\theta) = G_N(\theta)\cdots G_2(\theta)G_1(\theta).
\]

The quality measure used by the code is process fidelity:
\[
F(U_{\mathrm{target}}, U(x,\theta)) =
\frac{\left|\mathrm{Tr}\left(U_{\mathrm{target}}^\dagger U(x,\theta)\right)\right|^2}{d^2}.
\]

The baseline unitary-synthesis objective is therefore
\[
\min_{x,\theta} C_F(x,\theta)
=
1 - F(U_{\mathrm{target}}, U(x,\theta)).
\]

When Pareto mode is enabled, the training pipeline does not directly optimize only $C_F$. Instead, after warmup it uses a scalarized rollout cost
\[
C_i =
w_F (1 - F_i)
    + \mathbf{1}[F_i \ge \tau]
      \left(
        w_d \hat d_i + w_c \hat c_i
      \right),
\]
where:

\begin{itemize}[leftmargin=*]
\item $F_i$ is the fidelity of circuit $i$,
\item $d_i$ is its estimated depth,
\item $c_i$ is its CNOT count,
\item $\hat d_i$ and $\hat c_i$ are rollout-batch z-scores,
\item $\tau$ is the configured fidelity threshold,
\item $w = (w_F, w_d, w_c)$ is a Dirichlet-sampled weight vector with a hard lower bound on $w_F$.
\end{itemize}

So the current implementation should be viewed as solving
\[
\min_{x,\theta} 1 - F(U_{\mathrm{target}}, U(x,\theta))
\]
as the core synthesis objective, while optionally training toward a Pareto front over fidelity, depth, and CNOT count.

\section{Method Overview}

At a high level, the code follows this pipeline:

\begin{enumerate}[leftmargin=*]
\item Load and validate the experiment configuration.
\item Build the operator pool that defines the gate vocabulary.
\item Generate or load the target unitary.
\item Verify target unitarity and compile the target with Qiskit for reference.
\item Initialize a GPT-2 style autoregressive policy over gate tokens.
\item Sample batches of fixed-length gate sequences.
\item Evaluate sampled circuits with either pure discrete fidelity or hybrid discrete-continuous fidelity.
\item Store sampled sequences, scalar costs, and behavior log-probabilities in a replay buffer.
\item Train the policy with a clipped GRPO objective using replay-buffered data.
\item Maintain an optional Pareto archive of non-dominated circuits.
\item Optionally run a post-training Pareto gradient-descent refinement pass.
\item Decode the best circuit, verify it again, rebuild it as a Qiskit circuit, and compare it against the Qiskit reference compilation.
\end{enumerate}

This is therefore a hybrid search pipeline combining sequence modeling for structure search, matrix-based fidelity evaluation, gradient-based angle refinement, replay-buffered policy optimization, and optional multi-objective Pareto tracking.

\section{Current Default Configuration}

The current \texttt{config.yml} describes the following default experiment:

\begin{center}
\begin{tabular}{ll}
\toprule
Component & Current value \\
\midrule
Target qubits & \texttt{2} \\
Target type & \texttt{file} \\
Target path & \texttt{targets/haar\_random\_q2.npy} \\
Model size & \texttt{small} \\
Max generated gates & \texttt{28} \\
Training epochs & \texttt{200} \\
Rollout samples per epoch & \texttt{200} \\
Batch size & \texttt{32} \\
Learning rate & \texttt{1e-4} \\
Scheduler & \texttt{fixed} \\
Initial inverse temperature & \texttt{0.5} \\
Scheduler delta & \texttt{0.02} \\
Replay buffer size & \texttt{500} \\
Replay warmup size & \texttt{100} \\
Continuous optimization & enabled \\
Continuous optimizer & \texttt{lbfgs} \\
Continuous optimizer steps & \texttt{200} \\
Continuous optimizer restarts & \texttt{3} \\
\texttt{top\_k} filtering & \texttt{0}, meaning optimize all samples \\
Pareto archive & enabled \\
Pareto post-training GD & enabled \\
\bottomrule
\end{tabular}
\end{center}

The default run therefore exercises the most complex path in the repository: file-based target loading, continuous angle optimization during rollout scoring, Pareto scalarization during training, and a post-training Pareto refinement pass.

\section{Detailed Methodology}

\subsection{Operator Pool and Search Space}

The operator pool is the discrete search space used by the policy. In the current implementation, the sampling vocabulary contains:

\begin{itemize}[leftmargin=*]
\item one \texttt{RZ\_q$i$} token for each qubit $i$,
\item one \texttt{SX\_q$i$} token for each qubit $i$,
\item one \texttt{CNOT\_q$i$\_q$j$} token for each ordered control-target pair with $i \neq j$.
\end{itemize}

For $n$ qubits, the current pool size is
\[
|V| = n + n + n(n - 1) = n^2 + n.
\]

Important implementation details:

\begin{itemize}[leftmargin=*]
\item Although \texttt{continuous\_optimizer.py} knows how to parse \texttt{RX} and \texttt{RY}, the operator pool currently exposes only \texttt{RZ}, \texttt{SX}, and \texttt{CNOT} at generation time.
\item Rotation gates are stored in the pool with a placeholder angle of $\pi/4$.
\item The stored matrices are full dense $2^n \times 2^n$ unitaries embedded into the whole system.
\end{itemize}

For the default $n=2$ case:

\begin{itemize}[leftmargin=*]
\item $|V| = 6$,
\item the model vocabulary size is $7$ after adding the BOS token.
\end{itemize}

\subsection{Target Unitary Preparation}

Target construction is centralized in \texttt{target.py}. The supported modes are:

\begin{itemize}[leftmargin=*]
\item \texttt{random\_reachable}: compose a random short circuit from the operator pool;
\item \texttt{haar\_random}: generate a Haar-random unitary using QR decomposition of a complex Gaussian matrix;
\item \texttt{file}: load a saved unitary from a \texttt{.npy} file and validate its dimension.
\end{itemize}

With the current configuration, the target path is \texttt{targets/haar\_random\_q2.npy}, so the target is loaded from disk. After loading, \texttt{main.py} verifies numerically that
\[
U_{\mathrm{target}} U_{\mathrm{target}}^\dagger \approx I_d.
\]

\subsection{Reference Compilation with Qiskit}

Before running GQE, the target unitary is compiled by Qiskit using basis gates \texttt{rz}, \texttt{sx}, and \texttt{cx} with \texttt{optimization\_level = 3}. This result is not used for training. It is used only as a reference baseline for depth, total gate count, and two-qubit gate count.

\subsection{Sequence Representation}

The current model uses fixed-length sequences. There is:

\begin{itemize}[leftmargin=*]
\item one beginning-of-sequence token \texttt{BOS} with id \texttt{0},
\item a gate-token offset of \texttt{+1},
\item no end-of-sequence token,
\item no variable-length decoding in the main policy.
\end{itemize}

Each candidate circuit is therefore represented as a sequence of length
\[
N + 1,
\]
where the extra element is the BOS token.

For the current configuration:

\begin{itemize}[leftmargin=*]
\item $N = 28$,
\item every sampled candidate has exactly $29$ tokens including BOS,
\item every sampled candidate contains exactly $28$ generated gate tokens before any later simplification.
\end{itemize}

\subsection{Policy Model}

The structure-search model is implemented in \texttt{model.py} as a GPT-2 style causal transformer in Flax. It uses learned token embeddings, learned positional embeddings, causal self-attention, pre-layer-normalized transformer blocks, a GELU MLP, and tied output projection.

The supported presets are:

\begin{center}
\begin{tabular}{lrrr}
\toprule
Size & Layers & Heads & Embedding size \\
\midrule
\texttt{tiny} & 2 & 2 & 128 \\
\texttt{small} & 6 & 6 & 384 \\
\texttt{medium} & 12 & 12 & 768 \\
\texttt{large} & 24 & 16 & 1024 \\
\bottomrule
\end{tabular}
\end{center}

The default run uses \texttt{small}, so the current default policy has 6 transformer blocks, 6 attention heads, and embedding dimension 384.

One subtle but important detail is that the code treats the model outputs as energies rather than conventional language-model logits. Sampling and log-probability evaluation use \texttt{softmax(-beta * output)} instead of \texttt{softmax(output)}.

\subsection{Rollout Generation}

Rollouts are generated autoregressively in \texttt{Pipeline.generate()} with a JIT-compiled \texttt{jax.lax.scan}. At each generation step:

\begin{enumerate}[leftmargin=*]
\item the transformer receives the current token prefix,
\item the BOS token is masked out so it cannot be sampled again,
\item the next token is sampled with \texttt{jax.random.categorical} from \texttt{-beta * logits},
\item the sampled token is appended to the sequence.
\end{enumerate}

Formally, if $a_t$ is the token chosen at step $t$, the policy used by the code is
\[
\pi_\beta(a_t = v \mid a_{<t})
\propto
\exp\left(-\beta E_t(v)\right),
\]
where $E_t(v)$ is the model output for token $v$ at position $t$.

The sequence log-probability used later by GRPO is
\[
\log \pi_\beta(x)
=
\sum_{t=1}^{N}
\log \pi_\beta(x_t \mid x_{<t}).
\]

\subsection{Fidelity Evaluation}

The fidelity computation is defined in \texttt{cost.py}. Given a composed circuit unitary $U_{\mathrm{circuit}}$, the code evaluates
\[
F(U_{\mathrm{target}}, U_{\mathrm{circuit}})
=
\frac{\left|\mathrm{Tr}(U_{\mathrm{target}}^\dagger U_{\mathrm{circuit}})\right|^2}{d^2},
\]
and converts it into the cost
\[
C_F = 1 - F.
\]

The JAX implementation uses \texttt{jax.lax.scan} to compose gate matrices across positions, \texttt{jax.vmap} to score whole batches, and \texttt{jax.jit} to compile the evaluation path.

\subsection{Discrete and Hybrid Scoring Paths}

\paragraph{Pure discrete mode.}
If \texttt{continuous\_opt.enabled = false}, the code decodes token ids into operator-pool indices, composes the stored pool matrices, computes process fidelity directly, and returns $1 - F$. In this mode, parameterized rotations are evaluated at the placeholder angle $\pi/4$.

\paragraph{Hybrid discrete-continuous mode.}
If \texttt{continuous\_opt.enabled = true}, the code evaluates the circuit structure and then refines its continuous parameters before assigning the final fidelity.

Two sub-modes exist:

\begin{itemize}[leftmargin=*]
\item if \texttt{top\_k = 0}, every sampled circuit is continuously optimized;
\item if \texttt{top\_k > 0}, the pipeline first computes a cheap discrete score for all circuits, then continuously optimizes only the top $k$ structures.
\end{itemize}

With the current configuration, \texttt{top\_k = 0}, so every sampled circuit in every rollout is angle-optimized.

\subsection{Continuous Angle Optimization}

The continuous optimizer is implemented in \texttt{continuous\_optimizer.py}. For a fixed structure $x$, it solves
\[
\min_\theta 1 - F(U_{\mathrm{target}}, U(x,\theta)).
\]

The optimizer works by:

\begin{enumerate}[leftmargin=*]
\item parsing gate names into \texttt{GateSpec} objects,
\item encoding each circuit into fixed-shape arrays of gate type, qubit index, CNOT pair index, and padded angles,
\item building the circuit unitary differentiably in JAX,
\item optimizing the angles with either L-BFGS or Adam,
\item simplifying the circuit algebraically,
\item re-optimizing once if simplification produced a shorter parametric circuit.
\end{enumerate}

The default run uses L-BFGS with 200 steps and 3 restarts. The restart strategy is:

\begin{itemize}[leftmargin=*]
\item one start from the default initialization,
\item remaining starts sampled uniformly from $[-\pi,\pi]$ on the active parameter entries.
\end{itemize}

The simplification rules currently implemented are:

\begin{itemize}[leftmargin=*]
\item drop near-zero parameterized rotations,
\item merge consecutive same-axis rotations on the same qubit,
\item cancel adjacent identical CNOT gates.
\end{itemize}

Another important implementation detail is precision. The global JAX runtime enables x64, but \texttt{Factory.create\_continuous\_optimizer()} constructs the training-time optimizer with \texttt{fast\_runtime=True}, so the rollout-time continuous optimization uses float32/complex64 internally for speed.

\subsection{Structure Metrics}

The code tracks:

\begin{itemize}[leftmargin=*]
\item estimated depth,
\item total gate count,
\item CNOT count.
\end{itemize}

Depth is computed by a greedy per-qubit accumulation rule. For each token:

\begin{itemize}[leftmargin=*]
\item a single-qubit gate increments the depth counter of its qubit,
\item a two-qubit gate sets both qubit depths to $\max(\mathrm{depth}_{q_0}, \mathrm{depth}_{q_1}) + 1$.
\end{itemize}

This is an inexpensive structural proxy, not a full scheduling or transpilation pass.

\subsection{Pareto Scalarization and Archive}

If Pareto mode is enabled, the pipeline maintains a non-dominated archive over:

\begin{itemize}[leftmargin=*]
\item maximize fidelity,
\item minimize depth,
\item minimize CNOT count.
\end{itemize}

A circuit $a$ dominates circuit $b$ if:

\begin{itemize}[leftmargin=*]
\item $a.\mathrm{fidelity} \ge b.\mathrm{fidelity}$,
\item $a.\mathrm{depth} \le b.\mathrm{depth}$,
\item $a.\mathrm{cnot\_count} \le b.\mathrm{cnot\_count}$,
\item and at least one inequality is strict.
\end{itemize}

During training:

\begin{itemize}[leftmargin=*]
\item warmup uses fidelity-only cost,
\item after warmup, the pipeline samples a weight vector from a Dirichlet distribution,
\item it enforces a minimum fidelity weight $w_F^{\min}$,
\item it applies complexity penalties only when fidelity exceeds the threshold $\tau$.
\end{itemize}

The scalarized cost used for replay-buffer training is
\[
C_i =
w_F (1 - F_i)
    + \mathbf{1}[F_i \ge \tau]
      \left(
        w_d \frac{d_i - \mu_d}{\sigma_d + \varepsilon}
        +
        w_c \frac{c_i - \mu_c}{\sigma_c + \varepsilon}
      \right).
\]

The archive rejects points below a fidelity floor, removes dominated points, prunes by crowding distance if necessary, and reports a 2-D hypervolume over fidelity and CNOT count. With the current config, the fidelity floor starts at 0.5 and is raised to 0.9 at epoch 100.

\subsection{Replay Buffer and Training Data Reuse}

The policy is not trained purely on the latest rollout. Instead, sampled sequences are inserted into a FIFO replay buffer. Each buffer item stores:

\begin{itemize}[leftmargin=*]
\item the token sequence,
\item the scalar cost used for learning,
\item the sequence log-probability under the behavior policy that produced it.
\end{itemize}

Warmup is implemented as:

\begin{verbatim}
while len(buffer) < warmup_size:
    collect_rollout()
\end{verbatim}

Given the current configuration, \texttt{num\_samples = 200} and \texttt{warmup\_size = 100}, so one rollout is already enough to fill the warmup requirement. \texttt{BufferDataset} then exposes replay-buffer data for multiple passes per epoch via \texttt{steps\_per\_epoch}.

\subsection{GRPO Training Objective}

The actual policy update is implemented in \texttt{pipeline.py} as a clipped GRPO-style surrogate.

For a mini-batch of replayed trajectories:

\begin{itemize}[leftmargin=*]
\item let $C_i$ be the stored scalar cost,
\item let $A_i$ be the normalized advantage,
\item let $L_i$ be the current sequence log-probability,
\item let $L_i^{\mathrm{old}}$ be the stored behavior sequence log-probability,
\item let $r_i = \exp(L_i - L_i^{\mathrm{old}})$.
\end{itemize}

The code computes
\[
A_i =
\frac{\mathrm{mean}(C) - C_i}{\mathrm{std}(C) + \varepsilon},
\]
so lower-than-average cost yields positive advantage.

The clipped surrogate loss is
\[
\mathcal{L}_{\mathrm{GRPO}}
=
- \frac{1}{B}
  \sum_{i=1}^{B}
  \min\left(
    r_i A_i,\;
    \mathrm{clip}(r_i, 1-\epsilon, 1+\epsilon) A_i
  \right),
\]
where $\epsilon = \texttt{grpo\_clip\_ratio}$.

The optimizer on policy parameters is global-norm clipping followed by \texttt{optax.adamw}.

\subsection{Scheduler Behavior}

The scheduler provides the inverse temperature $\beta$ used during rollout generation and log-probability evaluation.

The code supports:

\begin{itemize}[leftmargin=*]
\item \texttt{fixed},
\item \texttt{cosine}.
\end{itemize}

The current config selects \texttt{fixed}, but in the current implementation \texttt{DefaultScheduler} still updates
\[
\texttt{current\_temperature} \leftarrow \texttt{current\_temperature} + \texttt{delta}
\]
after each rollout, followed by clamping to the configured interval.

\subsection{Best-Circuit Tracking and Final Decoding}

During training, the code tracks the best circuit from the most recent rollout batch using fidelity-only costs. After training:

\begin{enumerate}[leftmargin=*]
\item the best token sequence is decoded back into gate names;
\item if continuous optimization is enabled, the code re-runs the continuous optimizer on that structure for verification;
\item the optimized structure is converted into a Qiskit circuit;
\item the resulting circuit is compared against the Qiskit transpilation baseline.
\end{enumerate}

This final comparison reports fidelity, depth, total gates, and two-qubit gates.

\subsection{Post-Training Pareto Gradient Descent}

If \texttt{pareto\_gd.enabled = true}, the code performs a post-training pass over all archived Pareto circuits. For each archived circuit below $1 - \texttt{fidelity\_eps}$:

\begin{enumerate}[leftmargin=*]
\item convert the stored token sequence back to gate names,
\item run continuous optimization again,
\item update the archived fidelity if the optimizer improved it,
\item rebuild the archive from the updated points.
\end{enumerate}

This is intended to improve the final Pareto front without changing the learned policy itself.

\section{End-to-End Algorithm Summary}

The current implementation can be summarized as:

\begin{verbatim}
Input:
    config, target specification, operator pool

Build:
    U_target
    Qiskit reference circuit
    transformer policy
    replay buffer
    optional continuous optimizer
    optional Pareto archive

Warmup:
    sample rollouts
    compute fidelity-based costs
    fill replay buffer

Training epoch:
    generate num_samples fixed-length circuits
    compute old sequence log-probabilities
    evaluate fidelity, depth, and CNOT count
    if Pareto enabled and not in warmup:
        sample Pareto weights
        scalarize fidelity/depth/CNOT cost
    push sequences and costs to replay buffer
    train with clipped GRPO on replay-buffer batches
    update best-so-far statistics
    update scheduler
    update Pareto archive

After training:
    optionally run ParetoGD on archive members
    decode best structure
    verify best structure with continuous optimization
    build final Qiskit circuit
    compare against Qiskit baseline
\end{verbatim}

\section{Current Issues, Bugs, and Bottlenecks}

\subsection{Fixed-Length Generation Is a Hard Structural Limitation}

The policy has no EOS token and no variable-length decoding. Every sample contains exactly \texttt{max\_gates\_count} generated gates before any post hoc simplification.

Consequences:

\begin{itemize}[leftmargin=*]
\item the model cannot directly learn to stop early;
\item circuit-shortening pressure is only indirect;
\item the structural search space is larger than necessary;
\item reported complexity is partly shaped by simplification rather than native generation.
\end{itemize}

\subsection{\texttt{ParetoGDOptimizer} Does Not Write Back Simplified Structures}

The post-training Pareto GD pass calls \texttt{optimize\_circuit\_with\_params()}, but it only keeps the improved fidelity. It discards the returned optimized gate structure and parameters.

Consequences:

\begin{itemize}[leftmargin=*]
\item \texttt{token\_sequence}, \texttt{depth}, \texttt{total\_gates}, and \texttt{cnot\_count} in the archive remain unchanged;
\item if simplification found a shorter equivalent circuit, the archive does not reflect it;
\item the rebuilt archive can still be structurally stale after Pareto GD.
\end{itemize}

\subsection{Pareto Weight Sampling Is Not Fully Reproducible}

\texttt{Pipeline.\_sample\_weights()} uses \texttt{np.random.dirichlet(...)} from the global NumPy RNG instead of a seeded generator tied to \texttt{training.seed}.

\subsection{The ``Fixed'' Scheduler Is Not Actually Fixed}

The scheduler named \texttt{fixed} still increments its value by \texttt{delta} after every rollout, so experiment semantics are easy to misread.

\subsection{Dropout Exists in the Model but Is Disabled in Practice}

The transformer defines dropout layers, but the pipeline calls the model with \texttt{deterministic=True} during both generation and optimization, so no dropout regularization is active.

\subsection{Training-Time and Evaluation-Time Precision Differ}

The global runtime enables x64, but the training-time continuous optimizer uses float32/complex64 internally through \texttt{fast\_runtime=True}. The final verifier is not forced into the same low-precision path.

\subsection{Evaluation Cost Scales Poorly}

The current methodology uses full dense $2^n \times 2^n$ matrices for all operator-pool entries, circuit compositions, fidelity evaluations, and continuous optimization passes. With \texttt{top\_k = 0}, every sampled circuit in every rollout is angle-optimized.

\subsection{Depth Is a Proxy, Not a Hardware-Aware Metric}

Depth is estimated by a greedy qubit-depth counter on the raw token sequence. It is not computed from a simplified circuit, a scheduled circuit, or a transpiled hardware circuit.

\subsection{Logging and Naming Become Ambiguous in Pareto Mode}

Even when replay training uses scalarized multi-objective costs, \texttt{best\_cost} and \texttt{epoch\_best\_cost} are still tracked using fidelity-only costs.

\subsection{The Search Vocabulary Is Narrower Than the Optimizer's Capabilities}

The continuous optimizer understands \texttt{RX}, \texttt{RY}, and \texttt{RZ}, but the operator pool currently exposes only \texttt{RZ} as a parameterized single-qubit gate type.

\section{Suggestions for Improvement}

The most impactful next steps are:

\begin{enumerate}[leftmargin=*]
\item add variable-length generation with EOS and PAD handling;
\item fix Pareto GD so it updates structure as well as fidelity;
\item make Pareto runs fully reproducible by removing use of the global RNG;
\item separate scheduler names from scheduler behavior, or rename the current ``fixed'' schedule;
\item unify or explicitly manage the precision policy;
\item reduce the continuous-optimization bottleneck, for example by using \texttt{top\_k > 0} or a two-stage screening process;
\item compute complexity metrics from simplified or rebuilt circuits rather than raw token sequences;
\item expand the structural vocabulary, for example by enabling \texttt{RX} and \texttt{RY};
\item revisit replay-buffer strategy, including recency or Pareto-aware prioritization;
\item add regression tests around Pareto GD, reproducibility, scheduler semantics, and structure-metric consistency.
\end{enumerate}

\section{Closing Assessment}

The current codebase implements a coherent and technically interesting hybrid methodology for unitary synthesis:

\begin{itemize}[leftmargin=*]
\item a transformer proposes circuit structures,
\item a continuous optimizer repairs the continuous degrees of freedom,
\item GRPO updates the policy from replayed rollouts,
\item and Pareto machinery encourages better fidelity-complexity trade-offs.
\end{itemize}

Its strongest features are clear modular separation between target generation, structure search, and angle refinement; efficient JAX compilation of rollout and fidelity code paths; an explicit bridge to Qiskit for baseline comparison; and an emerging multi-objective training framework rather than fidelity-only search.

Its main current limitations are equally clear: fixed-length generation, expensive dense-matrix scaling, partial inconsistency between Pareto refinement and stored structure, and several experiment-control details that should be tightened for reproducibility and clarity.

That makes the project a strong small-scale research platform for unitary synthesis, with a clear path toward a more rigorous and scalable next version.

\end{document}
